\SourceLeaf{408}%
\[
\text{(2)}\quad .\ .\ .\ .\ .\quad\cos b=\cos a.\cos c+\sin a.\sin c.\cos B
\]
Dividing (A) by \(\sin c\), multiplying (2) by \(\operatorname{cotg}c\) and adding gives
\[
\frac{\cos a}{\sin c}=\frac{\cos a\cos c^2}{\sin c}+\sin b.\cos A+\sin a.\cos c.\cos B
\]
so
\[
\text{(3)}\quad .\ .\ .\quad\cos a.\sin c=\sin b.\cos A+\sin a.\cos c.\cos B
\]
and likewise
\[
\text{(4)}\quad .\ .\ .\quad\cos c.\sin a=\sin b.\cos C+\sin c.\cos a.\cos B
\]
\par
Multiplying (4) by \(\cos B\) and adding (3) gives
\[
\text{(5)}\quad .\ .\quad\cos a.\sin c.\sin B^2=\sin b.\cos A+\sin b.\cos B.\cos C
\]
and likewise
\[
\text{(6)}\quad .\ .\quad\cos a.\sin b.\sin C^2=\sin c.\cos A+\sin c.\cos B.\cos C
\]
\par
If we multiply (5) by \(\frac{\sin c}{\cos a}\) and subtract (6) multiplied by \(\frac{\sin b}{\cos a}\), we obtain
\[
\text{(7)}\quad .\ .\ .\ .\ .\quad\sin c^2.\sin B^2-\sin b^2.\sin C^2=0
\]
or
\[
\text{(B)}\quad .\ .\ .\ .\ .\ .\quad\sin c.\sin B=\sin b.\sin C
\]
Dividing (A,\,5) by \(\sin b\) and subtracting (B) multiplied by \(\frac{\cos a\sin B}{\sin b}\) gives
\[
\text{(C)}\quad .\ .\ .\ .\quad\cos A=-\cos B.\cos C+\sin B.\sin C.\cos a
\]
Finally, dividing (A,\,4) by \(\sin c\) and subtracting (B) multiplied by \(\frac{\operatorname{cotg}C}{\sin c}\) gives
\[
\text{(D)}\quad .\ .\ .\ .\quad\operatorname{cotg}c.\sin a-\operatorname{cotg}C.\sin B=\cos a.\cos B
\]
\par
From these four basic formulas follow at once the so-called analogies of \textsc{Napier} and the shortcuts available when the spherical triangle is required to be right-angled.

\{Compare this derivation with \textsc{Lagrange}'s in the sixth issue of the \textit{Journal de l’Ecole Polytechnique}.\}
